% GENERATED FILE. Edit canonical lessons, then run npm run generate:books.
% canonical-source-hash: fnv1a64:fa94bb44d96d214c
% canonical-lessons: GU-C160-pathari, GU-C160-patharvu, GU-C160-hincko, GU-C160-batti, GU-C160-holavvu

\chapter{Bedding, swing and light}
\label{ch:gu-bedding-swing-and-light}
\hlchaptermodality{160}

\begin{chapteropening}
\textbf{By the end of this chapter:} \emph{I can talk about making the bed and putting out the light.}

Complete the last lesson of chapter 160: I can talk about making the bed and putting out the light.
\end{chapteropening}

\section[\texorpdfstring{\gu{પથારી} (pathārī)}{પથારી (pathārī)}]{\gu{પથારી} (pathārī) — bedding, bed}

\label{lesson:GU-C160-pathari}



\begin{quote}
Before the new one: say the Gujarati for to spill, then the Gujarati for floor.
\end{quote}

\subsection*{You'll want to know: \gu{પથારી}}
\textbf{\gu{પથારી}} — \emph{pathārī} — \textquotedblleft{}bedding, bed\textquotedblright{}.

\begin{quote}
\textbf{\gu{મારી} \gu{પથારી} \gu{તૈયાર} \gu{છે}.} — \emph{mārī pathārī taiyār che} — \textquotedblleft{}My bed is ready.\textquotedblright{}
\end{quote}

\begin{cousinweb}
\textbf{\gu{પથારી}} (\emph{pathārī}) is a tadbhava. It comes from Sanskrit \emph{prastāra}, \textquotedblleft{}spreading out; bed of straw\textquotedblright{}, and CDIAL glosses the Gujarati word \textquotedblleft{}bedding\textquotedblright{} (CDIAL 8864).
\end{cousinweb}

\begin{grammarlens}[title={What you've built}]
One more word for everyday life at A2.
\end{grammarlens}

\subsection*{Your turn}
\begin{itemize}
\raggedright
  \item \emph{Say it:} \emph{pathārī}
  \item \emph{Say it:} \emph{pathārī}, once more
  \item \emph{Say it:} \emph{mārī pathārī taiyār che} — My bed is ready.
  \item \emph{Recall:} say the Gujarati for to spill, then the Gujarati for floor, then say \emph{pathārī} again
\end{itemize}

\begin{tcolorbox}[breakable,colback=teal!4,colframe=teal!35!black,title={Before you move on}]
What does \gu{પથારી} mean? (\textquotedblleft{}Bedding, bed\textquotedblright{}.) Say it once more.
\end{tcolorbox}
\section[\texorpdfstring{\gu{પાથરવું} (pātharvũ)}{પાથરવું (pātharvũ)}]{\gu{પાથરવું} (pātharvũ) — to spread out}

\label{lesson:GU-C160-patharvu}



\begin{quote}
Before the new one: say the Gujarati for floor, then the Gujarati for bedding.
\end{quote}

\subsection*{You'll want to know: \gu{પાથરવું}}
\textbf{\gu{પાથરવું}} — \emph{pātharvũ} — \textquotedblleft{}to spread out\textquotedblright{}.

\begin{quote}
\textbf{\gu{તે} \gu{રાત્રે} \gu{પથારી} \gu{પાથરે} \gu{છે}.} — \emph{te rātre pathārī pāthare che} — \textquotedblleft{}He spreads out the bedding at night.\textquotedblright{}
\end{quote}

\begin{cousinweb}
\textbf{\gu{પાથરવું}} (\emph{pātharvũ}) is a tadbhava, close kin of \emph{pathārī}. It comes from \emph{*prastarati}, \textquotedblleft{}spreads out\textquotedblright{}, a reconstructed form beside Sanskrit \emph{prastṛṇoti}; CDIAL glosses the Gujarati verb \textquotedblleft{}to spread\textquotedblright{} (CDIAL 8860).
\end{cousinweb}

\begin{grammarlens}[title={What you've built}]
One more word for everyday life at A2.
\end{grammarlens}

\subsection*{Your turn}
\begin{itemize}
\raggedright
  \item \emph{Say it:} \emph{pātharvũ}
  \item \emph{Say it:} \emph{pātharvũ}, once more
  \item \emph{Say it:} \emph{te rātre pathārī pāthare che} — He spreads out the bedding at night.
  \item \emph{Recall:} say the Gujarati for floor, then the Gujarati for bedding, then say \emph{pātharvũ} again
\end{itemize}

\begin{tcolorbox}[breakable,colback=teal!4,colframe=teal!35!black,title={Before you move on}]
What does \gu{પાથરવું} mean? (\textquotedblleft{}To spread out\textquotedblright{}.) Say it once more.
\end{tcolorbox}
\section[\texorpdfstring{\gu{હિંચકો} (hincko)}{હિંચકો (hincko)}]{\gu{હિંચકો} (hincko) — swing}

\label{lesson:GU-C160-hincko}



\begin{quote}
Before the new one: say the Gujarati for bedding, then the Gujarati for to spread out.
\end{quote}

\subsection*{You'll want to know: \gu{હિંચકો}}
\textbf{\gu{હિંચકો}} — \emph{hincko} — \textquotedblleft{}swing\textquotedblright{}.

\begin{quote}
\textbf{\gu{ઓસરીમાં} \gu{હિંચકો} \gu{છે}.} — \emph{osrīmā̃ hincko che} — \textquotedblleft{}There is a swing on the veranda.\textquotedblright{}
\end{quote}

\begin{cousinweb}
\textbf{\gu{હિંચકો}} (\emph{hincko}) is built on the inherited verb \emph{hĩcvũ}, \textquotedblleft{}to swing\textquotedblright{}, which CDIAL lists with \emph{hicakvũ} under \emph{*hicc}, \textquotedblleft{}draw back\textquotedblright{}, a reconstructed form (CDIAL 14083).
\end{cousinweb}

\begin{grammarlens}[title={What you've built}]
One more word for everyday life at A2.
\end{grammarlens}

\subsection*{Your turn}
\begin{itemize}
\raggedright
  \item \emph{Say it:} \emph{hincko}
  \item \emph{Say it:} \emph{hincko}, once more
  \item \emph{Say it:} \emph{osrīmā̃ hincko che} — There is a swing on the veranda.
  \item \emph{Recall:} say the Gujarati for bedding, then the Gujarati for to spread out, then say \emph{hincko} again
\end{itemize}

\begin{tcolorbox}[breakable,colback=teal!4,colframe=teal!35!black,title={Before you move on}]
What does \gu{હિંચકો} mean? (\textquotedblleft{}Swing\textquotedblright{}.) Say it once more.
\end{tcolorbox}
\section[\texorpdfstring{\gu{બત્તી} (battī)}{બત્તી (battī)}]{\gu{બત્તી} (battī) — electric light, lamp}

\label{lesson:GU-C160-batti}



\begin{quote}
Before the new one: say the Gujarati for to spread out, then the Gujarati for swing.
\end{quote}

\subsection*{You'll want to know: \gu{બત્તી}}
\textbf{\gu{બત્તી}} — \emph{battī} — \textquotedblleft{}electric light, lamp\textquotedblright{}.

\begin{quote}
\textbf{\gu{આ} \gu{બત્તી} \gu{નવી} \gu{છે}.} — \emph{ā battī navī che} — \textquotedblleft{}This light is new.\textquotedblright{}
\end{quote}

\begin{cousinweb}
\textbf{\gu{બત્તી}} (\emph{battī}) is usually traced to Sanskrit \emph{varti}, \textquotedblleft{}wick\textquotedblright{}, and EtymDB derives it from that word, but its path is not settled. CDIAL lists the kindred Gujarati \emph{vāṭ}, \textquotedblleft{}lamp\textquotedblright{}, and \emph{vātī} under that Sanskrit word, but not \emph{battī} itself (CDIAL 11359).
\end{cousinweb}

\begin{grammarlens}[title={What you've built}]
One more word for everyday life at A2.
\end{grammarlens}

\subsection*{Your turn}
\begin{itemize}
\raggedright
  \item \emph{Say it:} \emph{battī}
  \item \emph{Say it:} \emph{battī}, once more
  \item \emph{Say it:} \emph{ā battī navī che} — This light is new.
  \item \emph{Recall:} say the Gujarati for to spread out, then the Gujarati for swing, then say \emph{battī} again
\end{itemize}

\begin{tcolorbox}[breakable,colback=teal!4,colframe=teal!35!black,title={Before you move on}]
What does \gu{બત્તી} mean? (\textquotedblleft{}Electric light, lamp\textquotedblright{}.) Say it once more.
\end{tcolorbox}
\section[\texorpdfstring{\gu{હોલવવું} (holavvũ)}{હોલવવું (holavvũ)}]{\gu{હોલવવું} (holavvũ) — to put out, to switch off}

\label{lesson:GU-C160-holavvu}



\begin{quote}
Before the new one: say the Gujarati for swing, then the Gujarati for electric light.
\end{quote}

\subsection*{You'll want to know: \gu{હોલવવું}}
\textbf{\gu{હોલવવું}} — \emph{holavvũ} — \textquotedblleft{}to put out, to switch off\textquotedblright{}.

\begin{quote}
\textbf{\gu{હું} \gu{રાત્રે} \gu{બત્તી} \gu{હોલવું} \gu{છું}.} — \emph{hũ rātre battī holavũ chũ} — \textquotedblleft{}I switch off the light at night.\textquotedblright{}
\end{quote}

\begin{cousinweb}
\textbf{\gu{હોલવવું}} (\emph{holavvũ}) is a tadbhava. CDIAL lists it, glossed \textquotedblleft{}to extinguish\textquotedblright{}, under Sanskrit \emph{avaśrayaṇa}, \textquotedblleft{}taking a pot off the fire\textquotedblright{} (CDIAL 855a).
\end{cousinweb}

\begin{grammarlens}[title={What you've built}]
One more word for everyday life at A2.
\end{grammarlens}

\subsection*{Your turn}
\begin{itemize}
\raggedright
  \item \emph{Say it:} \emph{holavvũ}
  \item \emph{Say it:} \emph{holavvũ}, once more
  \item \emph{Say it:} \emph{hũ rātre battī holavũ chũ} — I switch off the light at night.
  \item \emph{Recall:} say the Gujarati for swing, then the Gujarati for electric light, then say \emph{holavvũ} again
\end{itemize}

\begin{tcolorbox}[breakable,colback=teal!4,colframe=teal!35!black,title={Before you move on}]
What does \gu{હોલવવું} mean? (\textquotedblleft{}To put out, to switch off\textquotedblright{}.) Say it once more.
\end{tcolorbox}
